% SPDX-License-Identifier: LPPL-1.3c
% Copyright (C) 2026 Christophe Jorssen
% Author and Current Maintainer: Christophe Jorssen <christophe.jorssen@gmail.com>
% LPPL maintenance status: maintained
% This file is part of luafemm. See LICENSE and LICENSES.md for its terms.
\documentclass{article}
\usepackage[margin=20mm]{geometry}
\usepackage{tikz,pgfplots}
\usetikzlibrary{femm,backgrounds}
\pgfplotsset{compat=1.18}
\pagestyle{empty}
\begin{document}
\usetikzlibrary{backgrounds}
\tikzset{
  lesson electric box/.style={femm/electrostatic={model=finite elements,
    xmin=-3,xmax=3,ymin=-2,ymax=2,mesh size=.3}},
  lesson electric plates/.pic={
    \draw[very thick,femm/electric electrode={name=positive,potential=50}]
      (-.5,-1)--(-.5,1);
    \draw[very thick,femm/electric electrode={name=negative,potential=-50}]
      (.5,-1)--(.5,1);}}

\par\medskip
\begin{tikzpicture}[lesson electric box,
    femm/electrostatic={name=fe coarse},scale=.8]
  \pic[transform shape] {lesson electric plates};
  \begin{scope}[on background layer]
    \pic {femm electric mesh};
    \pic[femm/electric={lines per charge=20}] {femm electric field};
  \end{scope}
  \draw[gray,densely dotted] (-3,-2) rectangle (3,2);
  \node[above] at (-.5,1) {$+50\,\mathrm V$};
  \node[below] at (.5,-1) {$-50\,\mathrm V$};
  \path[femm/electric profile={name=fe coarse axis,samples=101}]
    (0,-1.8)--(0,1.8);
\end{tikzpicture}

\par\medskip
\begin{tikzpicture}[lesson electric box,
    femm/electrostatic={mesh size=.15,name=fe fine},scale=.8]
  \pic[transform shape] {lesson electric plates};
  \begin{scope}[on background layer]
    \pic[femm/electric={lines per charge=24}] {femm electric field};
    \pic[femm/electric={levels={-40,-20,0,20,40}}] {femm equipotentials};
  \end{scope}
  \draw[gray,densely dotted] (-3,-2) rectangle (3,2);
  \path[femm/electric profile={name=fe fine axis,samples=161}]
    (0,-1.8)--(0,1.8);
\end{tikzpicture}

\par\medskip
\begin{tikzpicture}[lesson electric box,femm/electrostatic={
    xmin=-5,xmax=5,ymin=-5,ymax=5,mesh size=.15,name=fe wide}]
  \pic[transform shape] {lesson electric plates};
  \path[femm/electric seeds={samples=17}] (0,-1.8)--(0,1.8);
  \pic[femm/electric={xmin=-2,xmax=2,ymin=-2,ymax=2}]
    {femm electric field};
  \draw[gray,densely dotted] (-2,-2) rectangle (2,2);
  \path[femm/electric profile={name=fe wide axis,samples=161}]
    (0,-1.8)--(0,1.8);
\end{tikzpicture}

\par\medskip
\begin{tikzpicture}
\begin{axis}[width=10cm,height=5cm,no markers,grid=major,
  xlabel={$y$ (cm)},ylabel={$E_x$ (V/m)},
  legend style={at={(.5,1.03)},anchor=south,legend columns=3}]
  \femmelectricplot[profile=fe coarse axis,component=Ex,abscissa=y,
    plot options={gray,thick}]
  \addlegendentry{Coarse}
  \femmelectricplot[profile=fe fine axis,component=Ex,abscissa=y,
    plot options={blue!70!black,thick}]
  \addlegendentry{Fine}
  \femmelectricplot[profile=fe wide axis,component=Ex,abscissa=y,
    plot options={red!70!black,thick}]
  \addlegendentry{Wider box}
\end{axis}
\end{tikzpicture}

\par\medskip
\begin{tabular}{lr}
Calculation & Positive electrode charge per length (C/m)\\
Coarse & \femmelectricvalue[model=fe coarse,electrode=positive,
  quantity=charge per length]\\
Fine & \femmelectricvalue[model=fe fine,electrode=positive,
  quantity=charge per length]\\
Wider box & \femmelectricvalue[model=fe wide,electrode=positive,
  quantity=charge per length]
\end{tabular}

\end{document}
